% \chapter{Cronograma}
% \label{cap:cronograma}

% A seguir, serão apresentadas as tarefas do plano de pesquisa do trabalho em questão:

% \begin{description}
 
% \item[Tarefa 1:] Estudo da biblioteca de otimização matemática \texttt{gurobipy}, da linguagem Python, a fim de integrá-la ao modelos utilizados para alcançar possíveis melhorias nos resultados obtidos nesse trabalho.

% \item[Tarefa 2:] Estudo e testes de modelos de \textit{Deep Learning} e de \gls{PINN}, para observar se tais modelos conseguem se adaptar melhor ao problema e aos dados em questão.

% \item[Tarefa 3:] Atualização dos códigos de treinamento, ajustando os hiperparâmetros dos modelos já utilizados, integrando os novos modelos estudados, bem como a biblioteca \texttt{gurobipy}. Tais modificações têm como objetivo a obtenção de resultados melhores aos apresentados nesse trabalho, desse modo, continuará sendo utilizada a técnica de validação cruzada (\textit{cross-validation}) e um busca exaustiva para os novos testes.

% \item[Tarefa 4:] Escrita do texto da dissertação.

% \item[Tarefa 5:] Escrita de artigo sobre os resultados.

% \item[Tarefa 6:] Defesa de dissertação.

% \end{description}

% \begin{quadro}[ht]
% \centering
% \Caption{Cronograma de execução do trabalho}
% \label{quadro:cronograma}
% \resizebox{\textwidth}{!}{
% \begin{tabular}{|>{\raggedright\arraybackslash}p{1.5cm}|*{9}{>{\centering\arraybackslash}p{1.1cm}|}}
% \hline
% \multirow{2}{*}{Etapas} & \multicolumn{9}{c|}{2026} \\
% \cline{2-10}
%  & Abr & Mai & Jun & Jul & Ago & Set & Out & Nov & Dez \\
% \hline
% Tarefa 1 & X & & & & & & & &  \\
% \hline
% Tarefa 2 & & X & & & & & & &  \\
% \hline
% Tarefa 3 & & & X & X & X & & & &  \\
% \hline
% Tarefa 4 & X & X & X & X & X & X & X & &  \\
% \hline
% Tarefa 5 & & & & & X & X & X & X &  \\
% \hline
% Tarefa 6 & & & & & & & & & X  \\
% \hline
% \end{tabular}
% }
% \end{quadro}